%==================================================================
% LAPIS MESIN — pengaturan paket, tata letak, dan format
%
% Berkas ini TIDAK PERLU DIUBAH oleh mahasiswa.
% Seluruh data pribadi dan pilihan gaya sitasi diatur di a0-identitas.tex
%==================================================================

%------------------------------------------------------------------
% 1. BAHASA
% Paket babel bahasa Indonesia diganti pengaturan istilah manual agar
% template dapat dikompilasi di semua lingkungan, termasuk yang belum
% memasang paket babel-indonesian. Karena hyphenation dimatikan di
% bawah, tidak ada fitur yang hilang.
% Jika memakai Overleaf/MiKTeX dan ingin mengaktifkan babel, hapus blok
% \renewcommand di bawah lalu aktifkan baris berikut:
% \usepackage[bahasa]{babel}
%------------------------------------------------------------------
\usepackage[utf8]{inputenc}
\usepackage[T1]{fontenc}
\usepackage{etoolbox}

% Istilah struktural mengikuti \bahasanaskah pada a0-identitas.tex.
% Halaman legal (persetujuan, pengesahan, pernyataan) tetap Bahasa
% Indonesia karena merupakan dokumen resmi institusi.
\ifdefstring{\bahasanaskah}{en}{%
    \renewcommand{\chaptername}{CHAPTER}%
    \renewcommand{\figurename}{Figure}%
    \renewcommand{\tablename}{Table}%
    \providecommand{\appendixname}{Appendix}%
    \renewcommand{\appendixname}{Appendix}%
    \renewcommand{\contentsname}{TABLE OF CONTENTS}%
    \renewcommand{\listfigurename}{LIST OF FIGURES}%
    \renewcommand{\listtablename}{LIST OF TABLES}%
    \renewcommand{\bibname}{BIBLIOGRAPHY}%
    \newcommand{\judullampiran}{LIST OF APPENDICES}%
    \newcommand{\katasumber}{Source}%
}{%
    \renewcommand{\chaptername}{BAB}%
    \renewcommand{\figurename}{Gambar}%
    \renewcommand{\tablename}{Tabel}%
    \providecommand{\appendixname}{Lampiran}%
    \renewcommand{\appendixname}{Lampiran}%
    \renewcommand{\contentsname}{DAFTAR ISI}%
    \renewcommand{\listfigurename}{DAFTAR GAMBAR}%
    \renewcommand{\listtablename}{DAFTAR TABEL}%
    \renewcommand{\bibname}{DAFTAR PUSTAKA}%
    \newcommand{\judullampiran}{DAFTAR LAMPIRAN}%
    \newcommand{\katasumber}{Sumber}%
}

%------------------------------------------------------------------
% JENIS NASKAH: ta (Tugas Akhir) atau proposal (Proposal Tugas Akhir).
% Disetel oleh berkas master (laporan.tex / proposal.tex) SEBELUM \input
% berkas ini. Bila belum disetel, dianggap "ta".
%------------------------------------------------------------------
\providecommand{\jenisnaskah}{ta}
\ifdefstring{\jenisnaskah}{proposal}{%
    \newcommand{\labeltipe}{Proposal Tugas Akhir}%
    \newcommand{\labeltype}{Undergraduate Thesis Proposal}%
    \newcommand{\submitline}{}%
    \newcommand{\submitlineen}{}%
}{%
    \newcommand{\labeltipe}{Tugas Akhir}%
    \newcommand{\labeltype}{Undergraduate Thesis}%
    \newcommand{\submitline}{diajukan untuk memenuhi syarat\\ mendapatkan gelar \gelar}%
    \newcommand{\submitlineen}{submitted in partial fulfillment of the requirements\\ for the degree of \gelaren}%
}

% Mematikan hyphenation (pemenggalan kata) sesuai kelaziman naskah Indonesia
\hyphenpenalty=10000
\exhyphenpenalty=10000
\tolerance=9999
\emergencystretch=3em

%------------------------------------------------------------------
% 2. HURUF DAN TATA LETAK HALAMAN
%------------------------------------------------------------------
\usepackage{mathptmx}                     % Times New Roman
\usepackage[a4paper,left=4cm,top=3cm,right=3cm,bottom=3cm]{geometry}
\usepackage{setspace}
\onehalfspacing                           % spasi 1,5
\setlength{\parindent}{1.27cm}            % indentasi awal paragraf
\usepackage{indentfirst}                  % paragraf pertama ikut menjorok

%------------------------------------------------------------------
% 3. FORMAT JUDUL BAB DAN SUBBAB
%------------------------------------------------------------------
\usepackage{titlesec}
\renewcommand{\chaptertitlename}{BAB}     % istilah Indonesia (setelah titlesec)

\titleformat{\chapter}[display]
  {\normalfont\bfseries\centering}
  {\Large\chaptertitlename\ \thechapter}
  {0pt}
  {\Large}
\titlespacing*{\chapter}{0pt}{-20pt}{30pt}

\titleformat{\section}{\normalfont\bfseries}{\thesection}{1em}{}
\titleformat{\subsection}{\normalfont\bfseries}{\thesubsection}{1em}{}
\titleformat{\subsubsection}{\normalfont\bfseries}{\thesubsubsection}{1em}{}

% Penomoran subbab: 1.1, 1.1.1, 1.1.1.1
\renewcommand{\thesection}{\arabic{chapter}.\arabic{section}}
\renewcommand{\thesubsection}{\thesection.\arabic{subsection}}
\renewcommand{\thesubsubsection}{\thesubsection.\arabic{subsubsection}}
\setcounter{secnumdepth}{3}
\setcounter{tocdepth}{3}

%------------------------------------------------------------------
% 4. NOMOR HALAMAN
%------------------------------------------------------------------
\usepackage{fancyhdr}
\setlength{\headheight}{14.5pt}            % ruang untuk nomor di kepala halaman
% Gaya halaman biasa: nomor di KANAN ATAS.
\fancypagestyle{fancy}{\fancyhf{}\fancyhead[R]{\thepage}\renewcommand{\headrulewidth}{0pt}}
% Halaman pembuka BAB memakai gaya "plain" secara otomatis (perintah
% \chapter). Di sana nomor diletakkan di BAWAH-TENGAH.
\fancypagestyle{plain}{\fancyhf{}\fancyfoot[C]{\thepage}\renewcommand{\headrulewidth}{0pt}}
\pagestyle{fancy}

% Pembagian aturan front matter vs isi utama, dikaitkan ke \pagenumbering:
%  - \pagenumbering{roman}  (front matter) -> gaya "plain": SEMUA halaman
%    front matter bernomor di BAWAH-TENGAH, seperti semula.
%  - \pagenumbering{arabic} (isi utama)    -> gaya "fancy": halaman biasa
%    di KANAN ATAS, sedangkan pembuka BAB tetap di bawah-tengah (plain).
% Halaman ber-\thispagestyle{empty} (sampul, formulir) tetap tanpa nomor.
\let\origpagenumberingsaved\pagenumbering
\renewcommand{\pagenumbering}[1]{%
  \origpagenumberingsaved{#1}%
  \ifstrequal{#1}{arabic}{\pagestyle{fancy}}{\pagestyle{plain}}%
}

%------------------------------------------------------------------
% 5. DAFTAR ISI, GAMBAR, TABEL, DAN LAMPIRAN
%------------------------------------------------------------------
\usepackage{tocloft}
\renewcommand{\cfttoctitlefont}{\hfill\Large\bfseries}
\renewcommand{\cftaftertoctitle}{\hfill}
\renewcommand{\cftloftitlefont}{\hfill\Large\bfseries}
\renewcommand{\cftafterloftitle}{\hfill}
\renewcommand{\cftlottitlefont}{\hfill\Large\bfseries}
\renewcommand{\cftafterlottitle}{\hfill}
\renewcommand{\cftchapfont}{\bfseries}
\renewcommand{\cftchappagefont}{\bfseries}
\renewcommand{\cftchapdotsep}{\cftdotsep}   % titik-titik pada entri bab

% Daftar Lampiran (tidak disediakan kelas report, dibuat dengan tocloft)
\newlistof{lampiran}{lap}{\judullampiran}
\renewcommand{\cftlaptitlefont}{\hfill\Large\bfseries}
\renewcommand{\cftafterlaptitle}{\hfill}
% Cara memakai: \daftarlampiran{Lampiran A  Instrumen Penelitian}
\newcommand{\daftarlampiran}[1]{\addcontentsline{lap}{lampiran}{#1}}

%------------------------------------------------------------------
% 6. GAMBAR
%------------------------------------------------------------------
\usepackage{graphicx}
\graphicspath{{gambar/}}
\usepackage{float}                        % penempatan [H]
\usepackage{subcaption}                   % gambar berdampingan (a)(b)
\usepackage{pdflscape}                    % halaman lanskap (ikut memutar PDF)
\usepackage{adjustbox}                    % menyesuaikan objek terlalu lebar

\usepackage[font=small,labelsep=space,justification=centering]{caption}
% Keterangan gambar di BAWAH gambar, keterangan tabel di ATAS tabel
\captionsetup[figure]{position=below}
\captionsetup[table]{position=above}
% Penomoran mengikuti bab: Gambar 2.1, Tabel 4.3
\renewcommand{\thefigure}{\arabic{chapter}.\arabic{figure}}
\renewcommand{\thetable}{\arabic{chapter}.\arabic{table}}

% Perintah bantu untuk mencantumkan sumber di bawah gambar/tabel
\newcommand{\sumber}[1]{%
  \begin{center}\vspace{-0.6em}\footnotesize \katasumber: #1\end{center}}

% Penanda DATA YANG BELUM DIISI.
% Menghasilkan teks merah tebal di dalam kurung siku sehingga mustahil
% terlewat saat naskah dibaca ulang. Ganti dengan angka sebenarnya,
% lalu hapus perintah ini.
% Pemakaian: \belumdiisi{N} atau \belumdiisi{skor SUS}
\newcommand{\belumdiisi}[1]{\textbf{\textcolor{red}{[#1]}}}

% Kop formulir resmi FR-UAJ.
% Pemakaian: \kopformulir{FR-UAJ-07-46/R0}
%
% Halaman formulir mewarisi margin buku (kiri 4 cm, atas 3 cm) yang
% membuat kop terdorong ke tengah. Ketiga panjang di bawah menariknya
% kembali ke sudut kiri atas agar benar-benar tampak sebagai kop, tanpa
% mengubah margin isi formulir. Sesuaikan bila perlu.
\newlength{\kopnaik}     \setlength{\kopnaik}{2.25cm}   % seberapa jauh ditarik ke atas
\newlength{\kopkiri}     \setlength{\kopkiri}{1.6cm}    % seberapa jauh digeser ke kiri
\newlength{\koplebar}    \setlength{\koplebar}{4.6cm}   % lebar gambar kop

\newcommand{\kopformulir}[1]{%
  \vspace*{-\kopnaik}%
  \noindent\hspace*{-\kopkiri}%
  \begin{minipage}{\dimexpr\textwidth+\kopkiri\relax}%
    \noindent
    \begin{minipage}[t]{\dimexpr\koplebar+2mm\relax}%
      \vspace{0pt}\includegraphics[width=\koplebar]{kop-uaj}%
    \end{minipage}%
    \hfill
    \begin{minipage}[t]{5cm}%
      \vspace{2pt}\raggedleft\textbf{#1}%
    \end{minipage}%
  \end{minipage}%
  \par\vspace{1.4em}}

% Baris isian formulir (tanpa garis bawah).
% Pemakaian: \isianform{NAMA}{Budi Santoso}
%
% Lebar kolom label 4,8 cm dipilih agar label terpanjang "JUDUL TUGAS
% AKHIR" (tebal) muat penuh dan tidak menabrak tanda titik dua. Semua
% titik dua pada formulir jadi rata satu garis.
\newcommand{\isianform}[2]{%
  \noindent\makebox[4.8cm][l]{#1}: #2\par}

% Baris lanjutan (tanpa label), untuk isian yang lebih dari satu baris.
\newcommand{\isianlanjut}[1]{%
  \noindent\makebox[4.8cm][l]{}\phantom{: }#1\par}

% Isian judul yang MEMBUNGKUS OTOMATIS dari \judulid / \judulen.
% Judul panjang tidak perlu dipecah manual: teks mengalir mengikuti lebar
% kotak 8,6 cm dan berpindah baris sendiri, sejajar dengan kolom label.
% Gaya mengikuti halaman Pengesahan dan Keaslian (teks membungkus biasa).
% Pemakaian: \isianjudul{JUDUL TUGAS AKHIR}{\judulid}
\newcommand{\isianjudul}[2]{%
  \noindent\makebox[4.8cm][l]{#1}: %
  \parbox[t]{8.6cm}{\raggedright\setstretch{1.5}#2\par}\par}

% Penanda gambar yang BELUM tersedia berkasnya.
% Menghasilkan kotak bergaris putus-putus yang mustahil terlewat,
% sekaligus tetap terdaftar pada Daftar Gambar dan dapat dirujuk \cref.
% Pemakaian: \gambarbelumada{Keterangan gambar}{fig:label}
% Setelah berkas gambar tersedia, ganti dengan \includegraphics biasa.
\usepackage{tikz}
\newcommand{\gambarbelumada}[2]{%
  \begin{figure}[H]
    \centering
    \begin{tikzpicture}
      \draw[dashed,line width=0.9pt,color=gray]
        (0,0) rectangle (0.72\linewidth,4.2cm);
      \node[align=center,text=gray] at (0.36\linewidth,2.1cm)
        {\footnotesize\ttfamily [ BERKAS GAMBAR BELUM TERSEDIA ]\\[2pt]
         \footnotesize\rmfamily #1};
    \end{tikzpicture}
    \caption{#1}
    \label{#2}
  \end{figure}}

%------------------------------------------------------------------
% 7. TABEL
%------------------------------------------------------------------
\usepackage{longtable}                    % tabel bersambung antar halaman
\usepackage{multirow}                     % sel gabung vertikal
\usepackage{tabularx}                     % kolom lebar otomatis (tipe X)
\usepackage{threeparttable}               % tabel dengan catatan kaki
\usepackage{booktabs}                     % garis tabel gaya ilmiah
\usepackage{array}
\renewcommand{\arraystretch}{1.15}

%------------------------------------------------------------------
% 8. PERSAMAAN
%------------------------------------------------------------------
\usepackage{amsmath}
\usepackage{amssymb}
\renewcommand{\theequation}{\arabic{chapter}.\arabic{equation}}

%------------------------------------------------------------------
% 9. KODE PROGRAM DAN PSEUDOCODE
%------------------------------------------------------------------
\usepackage[table]{xcolor}
\usepackage{listings}
\definecolor{kodeabu}{RGB}{120,128,138}
\definecolor{kodehijau}{RGB}{0,110,60}
\definecolor{kodebiru}{RGB}{20,60,160}
\definecolor{kodelatar}{RGB}{247,248,250}
\lstset{
    backgroundcolor=\color{kodelatar},
    basicstyle=\ttfamily\footnotesize,
    keywordstyle=\color{kodebiru}\bfseries,
    commentstyle=\color{kodeabu}\itshape,
    stringstyle=\color{kodehijau},
    numbers=left, numberstyle=\tiny\color{kodeabu}, numbersep=8pt,
    frame=single, rulecolor=\color{kodeabu},
    breaklines=true, breakatwhitespace=false,
    showstringspaces=false, tabsize=4,
    captionpos=b, xleftmargin=16pt, framexleftmargin=16pt,
    literate={-}{-}1 {á}{{\'a}}1 {é}{{\'e}}1 {í}{{\'i}}1 {ó}{{\'o}}1
}
\renewcommand{\lstlistingname}{Kode Program}
% Penomoran kode program mengikuti bab: Kode Program 4.1, 4.2, dst.
\AtBeginDocument{\counterwithin{lstlisting}{chapter}}

% CATATAN: agar tidak meluber ke tepi halaman, usahakan setiap baris
% kode tidak lebih dari sekitar 60 karakter.

\usepackage[ruled,linesnumbered,lined]{algorithm2e}
\SetKwInput{KwMasukan}{Masukan}
\SetKwInput{KwKeluaran}{Keluaran}
\SetAlgorithmName{Algoritma}{Algoritma}{Daftar Algoritma}
% Penomoran algoritma mengikuti bab: Algoritma 3.1, 3.2, dst.
\AtBeginDocument{\counterwithin{algocf}{chapter}}

%------------------------------------------------------------------
% 10. DAFTAR (BULLET DAN NOMOR) RAPAT
%------------------------------------------------------------------
\usepackage{enumitem}
\newenvironment{daftarnomor}
  {\begin{enumerate}[topsep=2pt,itemsep=1pt,parsep=0pt,leftmargin=*]}
  {\end{enumerate}}
\newenvironment{daftarpoin}
  {\begin{itemize}[topsep=2pt,itemsep=1pt,parsep=0pt,leftmargin=*]}
  {\end{itemize}}

%------------------------------------------------------------------
% 11. SITASI DAN DAFTAR PUSTAKA
% Gaya diatur lewat \gayasitasi pada a0-identitas.tex
%------------------------------------------------------------------
\usepackage[round,authoryear]{natbib}
\bibpunct{(}{)}{;}{a}{,}{,}
\setlength{\bibhang}{1.27cm}              % indentasi gantung daftar pustaka
\setlength{\bibsep}{6pt}
\renewcommand{\bibsection}{}              % judul diatur manual di xx-daftarpustaka
\bibliographystyle{\gayasitasi}

%------------------------------------------------------------------
% 12. PRANALA DAN RUJUKAN SILANG
%------------------------------------------------------------------
\usepackage[hidelinks,bookmarks=true]{hyperref}
\usepackage{cleveref}
\crefname{figure}{Gambar}{Gambar}
\Crefname{figure}{Gambar}{Gambar}
\crefname{table}{Tabel}{Tabel}
\Crefname{table}{Tabel}{Tabel}
\crefname{equation}{Persamaan}{Persamaan}
\Crefname{equation}{Persamaan}{Persamaan}
\crefname{lstlisting}{Kode Program}{Kode Program}
\Crefname{lstlisting}{Kode Program}{Kode Program}
\crefname{algocf}{Algoritma}{Algoritma}
\Crefname{algocf}{Algoritma}{Algoritma}
\crefname{section}{Subbab}{Subbab}
\Crefname{section}{Subbab}{Subbab}
\crefname{chapter}{Bab}{Bab}
\Crefname{chapter}{Bab}{Bab}

% Kata sambung Bahasa Indonesia untuk rujukan jamak
\newcommand{\crefpairconjunction}{ dan }
\newcommand{\crefmiddleconjunction}{, }
\newcommand{\creflastconjunction}{, dan }
\newcommand{\crefrangeconjunction}{--}

% Rujukan silang ke rumusan masalah (RQ).
% Memakai \crefformat agar hasilnya "RQ1" tanpa spasi, bukan "RQ 1".
\crefformat{rq}{\labelrq#2#1#3}
\Crefformat{rq}{\labelrq#2#1#3}
\crefmultiformat{rq}{\labelrq#2#1#3}{ dan \labelrq#2#1#3}{, \labelrq#2#1#3}{, dan \labelrq#2#1#3}

% Bila naskah berbahasa Inggris, definisi di atas ditimpa di sini.
\ifdefstring{\bahasanaskah}{en}{%
    \crefname{figure}{Figure}{Figures}\Crefname{figure}{Figure}{Figures}%
    \crefname{table}{Table}{Tables}\Crefname{table}{Table}{Tables}%
    \crefname{equation}{Equation}{Equations}\Crefname{equation}{Equation}{Equations}%
    \crefname{section}{Section}{Sections}\Crefname{section}{Section}{Sections}%
    \crefname{chapter}{Chapter}{Chapters}\Crefname{chapter}{Chapter}{Chapters}%
    \crefname{lstlisting}{Code}{Codes}\Crefname{lstlisting}{Code}{Codes}%
    \crefname{algocf}{Algorithm}{Algorithms}\Crefname{algocf}{Algorithm}{Algorithms}%
    \renewcommand{\crefpairconjunction}{ and }%
    \renewcommand{\creflastconjunction}{, and }%
    \crefmultiformat{rq}{\labelrq#2#1#3}{ and \labelrq#2#1#3}{, \labelrq#2#1#3}{, and \labelrq#2#1#3}%
}{}
\crefname{tuj}{Tujuan}{Tujuan}
\Crefname{tuj}{Tujuan}{Tujuan}

% Penomoran otomatis rumusan masalah dan tujuan.
% Pemakaian:  \begin{rumusanmasalah} \item Pertanyaan\label{rq:kualitas} ... \end{rumusanmasalah}
% Merujuknya: \cref{rq:kualitas}  ->  menghasilkan "RQ1"
\newcounter{rq}
\renewcommand{\therq}{\arabic{rq}}
% Catatan penyetelan: agar label RQ sejajar dengan nomor subbab, posisi
% awal label harus nol. Pada lingkungan list, label mulai di
% (leftmargin - labelsep - labelwidth), sehingga ketiganya diatur agar
% leftmargin = labelsep + labelwidth. Tanpa ini label menjorok ke kiri
% sejauh \labelsep dan tidak rata dengan judul subbab.
\newenvironment{rumusanmasalah}
  {\setcounter{rq}{0}%
   \begin{list}{\labelrq\therq.}{\usecounter{rq}%
     \setlength{\labelwidth}{\lebarlabelrq}\setlength{\labelsep}{0.5em}%
     \setlength{\leftmargin}{\dimexpr\lebarlabelrq+0.5em\relax}%
     \setlength{\rightmargin}{0pt}\setlength{\listparindent}{0pt}%
     \setlength{\topsep}{3pt}\setlength{\itemsep}{2pt}\setlength{\parsep}{0pt}%
     \renewcommand{\makelabel}[1]{##1\hfill}}}
  {\end{list}}

\newcounter{tuj}
\renewcommand{\thetuj}{\arabic{tuj}}
\newenvironment{daftartujuan}
  {\setcounter{tuj}{0}%
   \begin{list}{\labeltujuan\thetuj.}{\usecounter{tuj}%
     \setlength{\labelwidth}{\lebarlabelrq}\setlength{\labelsep}{0.5em}%
     \setlength{\leftmargin}{\dimexpr\lebarlabelrq+0.5em\relax}%
     \setlength{\rightmargin}{0pt}\setlength{\listparindent}{0pt}%
     \setlength{\topsep}{3pt}\setlength{\itemsep}{2pt}\setlength{\parsep}{0pt}%
     \renewcommand{\makelabel}[1]{##1\hfill}}}
  {\end{list}}

%------------------------------------------------------------------
% 13. LAIN-LAIN
%------------------------------------------------------------------
\usepackage{lipsum}                       % teks contoh (boleh dihapus)
\usepackage{ragged2e}
\usepackage{lastpage}
